% GENERATED FILE. Edit canonical lessons, then run npm run generate:books.
% canonical-source-hash: fnv1a64:27f6639973a4d6f1
% canonical-lessons: LA-C417-possessio, LA-C417-proprius, LA-C417-alienus, LA-C417-dominium, LA-C417-mancipium

\chapter{Property}
\label{ch:la-property}
\hlchaptermodality{417}

\begin{chapteropening}
\textbf{By the end of this chapter:} \emph{I can talk about possession and what is mine or another's, ownership and formal purchase.}

Complete the last lesson of chapter 417: I can talk about possession and what is mine or another's, ownership and formal purchase.
\end{chapteropening}

\section[possessiō, possessiōnis]{possessiō, possessiōnis — possession; (plural) estates}

\label{lesson:LA-C417-possessio}



\begin{quote}
Before the new one: say the Latin for to restrain, then the Latin for to claim.
\end{quote}

\subsection*{You'll want to know: possessiō, possessiōnis}
\textbf{possessiō, possessiōnis} — \textquotedblleft{}possession; (plural) estates\textquotedblright{}.

\begin{quote}
Possessiōnēs magnās in agrīs habēbat.
\end{quote}

\begin{quote}
\emph{He had large estates in the country.}
\end{quote}

\emph{possessiō, possessiōnis} f., the noun of \emph{possideō}. In the plural, landed estates. English \emph{possession}.

\begin{cousinweb}
\emph{possessiō} is the noun built on \emph{possideō} \textquotedblleft{}to possess\textquotedblright{} (EtymDB; Lewis \& Short). English \emph{possession} comes from it.
\end{cousinweb}

\begin{grammarlens}[title={What you've built}]
One more word for telling, explaining and describing at B1.
\end{grammarlens}

\subsection*{Your turn}
\begin{itemize}
\raggedright
  \item \emph{Say it:} \emph{possessiō, possessiōnis}
  \item \emph{Say it:} \emph{possessiō, possessiōnis}, once more
  \item \emph{Say it:} a sentence of your own with \emph{possessiō, possessiōnis}
  \item \emph{Recall:} say the Latin for to restrain, then the Latin for to claim, then say \emph{possessiō, possessiōnis} again
\end{itemize}

\begin{tcolorbox}[breakable,colback=teal!4,colframe=teal!35!black,title={Before you move on}]
What does possessiō, possessiōnis mean? (\textquotedblleft{}Possession; (plural) estates\textquotedblright{}.) Say it once more.
\end{tcolorbox}
\section[proprius, propria, proprium]{proprius, propria, proprium — one's own; peculiar (to)}

\label{lesson:LA-C417-proprius}



\begin{quote}
Before the new one: say the Latin for to claim, then the Latin for possession.
\end{quote}

\subsection*{You'll want to know: proprius, propria, proprium}
\textbf{proprius, propria, proprium} — \textquotedblleft{}one's own; peculiar (to)\textquotedblright{}.

\begin{quote}
Domum propriam nōn habet.
\end{quote}

\begin{quote}
\emph{He has no house of his own.}
\end{quote}

\emph{proprius, propria, proprium} + genitive or dative: belonging to one alone. English \emph{proper} and \emph{property}.

\begin{cousinweb}
The origin of \emph{proprius} is not settled: EtymDB traces it to Proto-Italic \emph{*prijos}, but Lewis \& Short call it doubtful and suggest, with a \textquotedblleft{}perhaps\textquotedblright{}, the root of \emph{prae} \textquotedblleft{}before\textquotedblright{}. English \emph{proper} and \emph{property} come from it.
\end{cousinweb}

\begin{grammarlens}[title={What you've built}]
One more word for telling, explaining and describing at B1.
\end{grammarlens}

\subsection*{Your turn}
\begin{itemize}
\raggedright
  \item \emph{Say it:} \emph{proprius, propria, proprium}
  \item \emph{Say it:} \emph{proprius, propria, proprium}, once more
  \item \emph{Say it:} a sentence of your own with \emph{proprius, propria, proprium}
  \item \emph{Recall:} say the Latin for to claim, then the Latin for possession, then say \emph{proprius, propria, proprium} again
\end{itemize}

\begin{tcolorbox}[breakable,colback=teal!4,colframe=teal!35!black,title={Before you move on}]
What does proprius, propria, proprium mean? (\textquotedblleft{}One's own; peculiar (to)\textquotedblright{}.) Say it once more.
\end{tcolorbox}
\section[aliēnus, aliēna, aliēnum]{aliēnus, aliēna, aliēnum — belonging to another; foreign; unfavourable}

\label{lesson:LA-C417-alienus}



\begin{quote}
Before the new one: say the Latin for possession, then the Latin for one's own.
\end{quote}

\subsection*{You'll want to know: aliēnus, aliēna, aliēnum}
\textbf{aliēnus, aliēna, aliēnum} — \textquotedblleft{}belonging to another; foreign; unfavourable\textquotedblright{}.

\begin{quote}
Rem aliēnam cupere nōn dēbēmus.
\end{quote}

\begin{quote}
\emph{We should not desire what belongs to another.}
\end{quote}

\emph{aliēnus, aliēna, aliēnum}, built on \emph{alius}, other. The opposite of \emph{suus} and \emph{proprius}. English \emph{alien}.

\begin{cousinweb}
\emph{aliēnus} is built on \emph{alius} \textquotedblleft{}other\textquotedblright{} (EtymDB; Lewis \& Short). English \emph{alien} comes from it through French.
\end{cousinweb}

\begin{grammarlens}[title={What you've built}]
One more word for telling, explaining and describing at B1.
\end{grammarlens}

\subsection*{Your turn}
\begin{itemize}
\raggedright
  \item \emph{Say it:} \emph{aliēnus, aliēna, aliēnum}
  \item \emph{Say it:} \emph{aliēnus, aliēna, aliēnum}, once more
  \item \emph{Say it:} a sentence of your own with \emph{aliēnus, aliēna, aliēnum}
  \item \emph{Recall:} say the Latin for possession, then the Latin for one's own, then say \emph{aliēnus, aliēna, aliēnum} again
\end{itemize}

\begin{tcolorbox}[breakable,colback=teal!4,colframe=teal!35!black,title={Before you move on}]
What does aliēnus, aliēna, aliēnum mean? (\textquotedblleft{}Belonging to another; foreign; unfavourable\textquotedblright{}.) Say it once more.
\end{tcolorbox}
\section[dominium, dominiī]{dominium, dominiī — ownership}

\label{lesson:LA-C417-dominium}



\begin{quote}
Before the new one: say the Latin for one's own, then the Latin for belonging to another.
\end{quote}

\subsection*{You'll want to know: dominium, dominiī}
\textbf{dominium, dominiī} — \textquotedblleft{}ownership\textquotedblright{}.

\begin{quote}
Dominium agrī ad fīlium pervēnit.
\end{quote}

\begin{quote}
\emph{The ownership of the field passed to the son.}
\end{quote}

\emph{dominium, dominiī} n., built on \emph{dominus}, master. In Roman law, full ownership. English \emph{dominion}.

\begin{cousinweb}
\emph{dominium} is built on \emph{dominus} \textquotedblleft{}a master\textquotedblright{} (EtymDB), which comes from \emph{domus} \textquotedblleft{}a house\textquotedblright{} (Lewis \& Short): the right of the master of the house. English \emph{dominion} is of the same family.
\end{cousinweb}

\begin{grammarlens}[title={What you've built}]
One more word for telling, explaining and describing at B1.
\end{grammarlens}

\subsection*{Your turn}
\begin{itemize}
\raggedright
  \item \emph{Say it:} \emph{dominium, dominiī}
  \item \emph{Say it:} \emph{dominium, dominiī}, once more
  \item \emph{Say it:} a sentence of your own with \emph{dominium, dominiī}
  \item \emph{Recall:} say the Latin for one's own, then the Latin for belonging to another, then say \emph{dominium, dominiī} again
\end{itemize}

\begin{tcolorbox}[breakable,colback=teal!4,colframe=teal!35!black,title={Before you move on}]
What does dominium, dominiī mean? (\textquotedblleft{}Ownership\textquotedblright{}.) Say it once more.
\end{tcolorbox}
\section[mancipium, mancipiī]{mancipium, mancipiī — formal purchase; a slave}

\label{lesson:LA-C417-mancipium}



\begin{quote}
Before the new one: say the Latin for belonging to another, then the Latin for ownership.
\end{quote}

\subsection*{You'll want to know: mancipium, mancipiī}
\textbf{mancipium, mancipiī} — \textquotedblleft{}formal purchase; a slave\textquotedblright{}.

\begin{quote}
Servum mancipiō accēpit.
\end{quote}

\begin{quote}
\emph{He received the slave by formal purchase.}
\end{quote}

\emph{mancipium, mancipiī} n., \emph{manus} + \emph{capiō}: taking in the hand, the old ceremony of purchase with scales and bronze; then the thing bought, a slave. English \emph{emancipate}.

\begin{cousinweb}
\emph{mancipium} is built on \emph{manceps} \textquotedblleft{}a purchaser\textquotedblright{} (EtymDB; Lewis \& Short), from \emph{manus} \textquotedblleft{}hand\textquotedblright{} and \emph{capiō} \textquotedblleft{}to take\textquotedblright{}. English \emph{emancipate} comes from its family.
\end{cousinweb}

\begin{grammarlens}[title={What you've built}]
One more word for telling, explaining and describing at B1.
\end{grammarlens}

\subsection*{Your turn}
\begin{itemize}
\raggedright
  \item \emph{Say it:} \emph{mancipium, mancipiī}
  \item \emph{Say it:} \emph{mancipium, mancipiī}, once more
  \item \emph{Say it:} a sentence of your own with \emph{mancipium, mancipiī}
  \item \emph{Recall:} say the Latin for belonging to another, then the Latin for ownership, then say \emph{mancipium, mancipiī} again
\end{itemize}

\begin{tcolorbox}[breakable,colback=teal!4,colframe=teal!35!black,title={Before you move on}]
What does mancipium, mancipiī mean? (\textquotedblleft{}Formal purchase; a slave\textquotedblright{}.) Say it once more.
\end{tcolorbox}
